\subsection{Every filling admits every even Euler number after vertical blowups}

\begin{statement}{C.2.UN.4}{Proved}{Every even Euler number occurs above each fixed torus filling}
Let $f:Y\to B$ be a proper holomorphic map from a smooth complex
threefold to a smooth complex curve, with connected fibres and smooth
general fibre a complex two-torus. Fix $b\in B$ and put
$S=(Y_b)_{\rm red}$. For every integer $k$ there is a finite composition
of blowups with smooth compact centres in successive fibres over $b$,
\[
 \pi:\widetilde Y\longrightarrow Y,\qquad
 \widetilde f=f\pi,\qquad
 e((\widetilde Y_b)_{\rm red})=2k.                 \tag{C.VR.1}
\]
The map $\pi$ is an isomorphism over $B\setminus\{b\}$.
In particular every such filling is dominated by a smooth filling
whose central reduced support has Euler number zero. More precisely,
every further smooth model obtained by vertical blowups admits a
zero-Euler refinement by additional vertical blowups. Thus zero-Euler
models are cofinal for this domination relation, in the explicit sense
that every object is dominated by one of them.
Above every given filling there are such zero-Euler refinements with
arbitrarily many irreducible components of their central support.

An elementary operation lowering the Euler number by two consists of
one point blowup and one smooth genus-three curve blowup. On integral
cohomology its additional summands are $\Z^2$ in degrees two and four
and $\Z^6$ in degree three. In particular, whenever the Betti numbers
are finite,
\[
 (\Delta b_2,\Delta b_3,\Delta b_4)=(2,6,2),
 \qquad \Delta b_j=0\quad(j\notin\{2,3,4\}).       \tag{C.VR.2}
\]
All the original positive full-fibre multiplicities are retained on
their strict transforms. The construction requires no normal crossings
condition on the original support.
\end{statement}
\begin{proof}
We first specify the comparison for a blowup whose centre meets
singularities of the support. Let $Z\subset S$ be a smooth compact
centre of complex codimension $r$ in the smooth ambient threefold,
and let $\pi_Z:Y'\to Y$ be its blowup. Write $S'=(f\pi_Z)^{-1}(b)_{\rm red}$.
The equality of underlying sets is $S'=\pi_Z^{-1}(S)$, and the actual
restrictions of the blowup give
\[
 S'\setminus\pi_Z^{-1}(Z)\simeq S\setminus Z,
 \qquad \pi_Z^{-1}(Z)=\PP(N_{Z/Y}).                \tag{C.VR.3}
\]
Compact complex analytic spaces admit finite triangulations, compatible
with these closed analytic subsets. Euler additivity, applied with
compact supports on the complements, and the projective-bundle
formula therefore give
\[
 \begin{split}
 e(S')-e(S)
 &=e(\PP(N_{Z/Y}))-e(Z)\\
 &=(r-1)e(Z).
 \end{split}                                      \tag{C.VR.4}
\]
The projective-bundle Euler equality also follows from its locally
trivial fibre $\PP^{r-1}$, whose Euler number is $r$.
Thus (C.VR.4) retains any intersections of the centre with other
components or with the singular support. It makes no replacement of
the full inverse image by a strict transform.

Choose a smooth point $p$ of one component of $S$, away from every
other component. Such points form a nonempty open subset of each
irreducible component. If the full fibre has multiplicity $m>0$ on
that component, local holomorphic coordinates may be chosen with
\[
                    t\circ f=x^m.                \tag{C.VR.5}
\]
Indeed the local equation initially has the form $x^m a$, where $a$
is a holomorphic unit; a local holomorphic $m$th root of $a$ replaces
$x$ by $x a^{1/m}$. Blow up $p$. The exceptional surface
$E\simeq\PP^2$ meets the strict component in the line
$L=\{X=0\}\subset\PP(T_pY)$, with homogeneous coordinates
$[X:Y:Z]$ induced by $(x,y,z)$. It meets no other central component.
In the chart $y=r$, $x=ru$, $z=rv$, the actual map is
\[
                         t=(ru)^m.                \tag{C.VR.6}
\]
Consequently both $E=(r=0)$ and the strict component $(u=0)$
have multiplicity $m$. Formula (C.VR.4), with $r=3$ and centre a
point, shows that this step increases the central Euler number by two.

Inside $E$ take the Fermat quartic
\[
 C=\{X^4+Y^4+Z^4=0\}\subset\PP^2.                \tag{C.VR.7}
\]
Its three partial derivatives cannot vanish simultaneously at a
projective point, so it is smooth. It is connected: if the quartic
were a disjoint union of distinct plane components, two such
components would meet by their positive intersection number; their
intersection would be a singular point. Adjunction gives
$K_C=\OO_C(1)$ and $2g(C)-2=4$, hence $g(C)=3$ and $e(C)=-4$.
On $L$ its equation is $Y^4+Z^4=0$. It has four distinct roots,
none at $Z=0$, and the two gradients of the quartic and $X$ are
independent there. Thus $C$ meets $L$ transversely in exactly four
points.

Now blow up the smooth curve $C$ in the threefold obtained after the
first blowup. Its exceptional divisor is the projective line bundle
$\PP(N_{C/Y_1})\to C$. Equation (C.VR.4), with $r=2$, gives Euler
change $e(C)=-4$. The two successive changes sum to $-2$.
For clarity, the four points of $C\cap L$ introduce no singularities
of the ambient space and do not alter this Euler comparison. Near any
one of them, preserve the local equations $r=0$ of $E$ and $u=0$ of
the strict old component. Transversality gives a third coordinate $v$
such that $C=(r=v=0)$. Absorbing a remaining unit in $r$ preserves
these equations and gives $t=(ru)^m$ as in (C.VR.6). The two charts of
the curve blowup are
\[
 \begin{array}{c|c|c}
 \text{chart}&\text{blowdown substitution}&\text{base map}\\ \hline
 \text{$r$ chart}&v=rw&t=(ru)^m\\
 \text{$v$ chart}&r=vw&t=(uvw)^m .
 \end{array}                                      \tag{C.VR.8}
\]
In the second chart the new exceptional divisor, the strict transform
of $E$, and the strict old component have equations $v=0$, $w=0$,
and $u=0$, respectively. They meet with normal crossings and all
have multiplicity $m$. Away from $C\cap L$ the comparison is the
same two-coordinate blowup of $t=r^m$. The strict transform of $E$
is isomorphic to $E$, since $C$ is Cartier in $E$; the strict old
component is its point blowup at the four transverse intersection
points. These identifications follow directly from the charts in
(C.VR.8). Outside the specified neighbourhoods the maps are
isomorphisms, so any other singularities of the original support
are simply retained.

We include the integral cohomology comparison giving (C.VR.2).
For a smooth centre of codimension $r$, let $j:E_Z\hookrightarrow Y'$
and $p:E_Z\to Z$ denote the exceptional inclusion and projection,
and put $\xi=c_1(\OO_{E_Z}(1))$. The actual blowup homomorphism is
\[
 \begin{split}
 H^q(Y,\Z)\oplus\bigoplus_{i=1}^{r-1}H^{q-2i}(Z,\Z)
   &\longrightarrow H^q(Y',\Z),\\
 (\alpha,\beta_1,\ldots,\beta_{r-1})
   &\longmapsto \pi_Z^*\alpha+
       \sum_{i=1}^{r-1}j_*(\xi^{i-1}p^*\beta_i).
 \end{split}                                      \tag{C.VR.9}
\]
It is an isomorphism. To check it integrally, excision identifies the
relative cohomology of $(Y',E_Z)$ with that of $(Y,Z)$. The
projective-bundle theorem makes $H^q(E_Z,\Z)$ the direct sum of
$\xi^i p^*H^{q-2i}(Z,\Z)$ for $0\le i<r$, and the self-intersection
formula gives $j^*j_*\eta=-\xi\eta$. Thus the extra terms of
(C.VR.9) map isomorphically, with sign minus, onto the summands with
$i>0$. Comparing the relative long exact sequences proves
surjectivity; their remaining $i=0$ term is exactly the pullback from
$Y$. The proper degree-one trace makes $\pi_Z^*$ injective, and
restriction to $E_Z$ then proves independence of all the displayed
extra summands. This proves (C.VR.9). The same argument applies on a
compact fibre tube with boundary, since the centre is in its interior
and the blowup is unchanged on the boundary.

The point blowup adds $\Z$ in each of degrees two and four. The
genus-three curve blowup adds $H^{q-2}(C,\Z)$, whose ranks in these
three successive degrees are $1,6,1$. This proves (C.VR.2) and shows
that pre-existing integral torsion is unchanged.

Both steps are proper holomorphic maps from smooth threefolds. A
blowup has connected fibres and pushes its structure sheaf to the
structure sheaf of the original smooth space. Stein factorization
therefore proves that their composition with $f$ again has connected
fibres. The centres lie over $b$, so the resulting map is the same
proper torus family over $B\setminus\{b\}$.

Finally C.2.UN.1 writes $e(S)=2n$ for an integer $n$. To obtain $2k$
with $k\ge n$, use $k-n$ point blowups. To obtain $2k$ with $k<n$,
use $n-k$ of the two-step operations just proved. At each stage a
fresh smooth point of a single component exists, and the construction
can be performed there; only finitely many centres are needed.
Applying this same procedure to any given vertical blowup model
proves the asserted cofinality with its specified domination maps.
After reaching Euler number zero, compose one further point blowup
with one of the two-step decreasing operations, at fresh smooth
points. Their net Euler change is zero, while they add three
irreducible central components. Repetition gives arbitrarily many
components above the original filling. Each repetition adds
$(3,6,3)$ to $(b_2,b_3,b_4)$ when those numbers are finite; the
actual maps remain the successive blowdowns already described.
The exact chart, quartic, Euler, and cohomology calculations are
retained in \path{code/E_vertical_refinements.py} and its output.
\end{proof}
